\section{考察}
\subsection{現在の構成との比較のまとめ}
表\ref{tab_summary}に，CLIPとソフトマックスを用いるhybridを基準とした3つの変更の効果をまとめる．
シグモイドへの変更は，深さ0の誤りを深い段階で検出できないという現在の構成の課題に対し，誤った枝の識別性能を改善する点で有効である．
一方，深さ0が正しい場合の深さ1の判断ではソフトマックスの方が識別性能が高く，停止規則の$ hF $の改善はhybridで0.012にとどまった．
両者の性質は相補的であり，各段階の語の選択と兄弟の語の間の比較にはソフトマックス，誤った枝の検出と棄却にはシグモイドを用いる組み合わせが考えられる．

\begin{table}[t]
    \centering
    \caption{CLIPとソフトマックスを用いるhybridを基準とした3つの変更の効果（テスト用の分割，YOLO11xの閾値0.25）．}
    \label{tab_summary}
    \small
    \setlength{\tabcolsep}{2pt}
    \begin{tabular}{lcc}
        \toprule
        変更 & 改善した指標 & 悪化した指標 \\
        \midrule
        シグモイド & \shortstack{誤枝AUROC \\ 0.418 → 0.617} & \shortstack{深さ1のAUROC \\ 0.762 → 0.644} \\
        \midrule
        スコアの比較 & -- & \shortstack{Top-1 0.969 → 0.936 \\（ソフトマックス）} \\
        \midrule
        \shortstack{WordNetの \\ 同じ深さの語彙} & \shortstack{COCO外の正解率 \\ 0 → 3.8\%} & \shortstack{Top-1 \\ 0.969 → 0.945} \\
        \bottomrule
    \end{tabular}
\end{table}

YOLOとVLMのスコアの比較は，尺度の異なる2つの値を直接比較するため，正規化関数を変えるだけで切替の頻度が6件から331件まで変化した．
2つのモデルのスコアを比較するためには，検証用の分割で各モデルの確信度を正解率に対応付ける較正を行った上で比較する必要がある．

WordNetの同じ深さの語彙は，COCOと同じ粒度の語でCOCO外の物体を命名するという目的には合致するが，現在の語彙では効果が小さい．
第一の理由は，表\ref{tab_wn_vocab}から，lvis\_testのCOCO外の物体のカテゴリの半数以上が語彙に上位語または自身を持たないことである．
第二の理由は，WordNetのレンマには物体の部分（leg，screen）や多義の語（slit，worsted）が含まれ，VLMがこれらを選ぶとCOCOの物体の正解が失われることである．
また，同じ深さの制約により，penguinやcamelのように直観的には基本的な語であってもCOCOのいとこより深い語は語彙から除かれる．
WordNetの深さと人が用いる名前の粒度は一致しないため，具体性の評定や画像データセットでの出現頻度等，物体の名前としての用いられやすさを考慮した語彙の選択が必要である．

\subsection{論文への反映}
論文の主張に対し，シグモイドの結果は，深い段階の確信度が誤りの手がかりにならないという課題に対する有効な対策として追記できる．
スコアの比較とWordNetの同じ深さの語彙は，現在の結果ではhybrid\_switch\_c0.5を上回らないため，提案手法として採用するのではなく，COCO外の語彙への体系的な拡張の検討結果と今後の課題として記述するのが妥当である．
